\section{جمع‌بندی}

در این پروژه، یک تحلیل داده‌کاوی بر روی مجموعه‌دادهٔ تجمیع‌شدهٔ بیماری قلبی از مخزن
\lr{UCI Machine Learning Repository} انجام شد که هدف آن پیش‌بینی وجود بیماری قلبی با استفاده از
طبقه‌بند \lr{K-Nearest Neighbors (KNN)} بود. مسیر انجام کار شامل پیش‌پردازش داده‌ها، تحلیل اکتشافی
جامع (\lr{EDA}) و پیاده‌سازی یک نسخهٔ ترکیبی از الگوریتم \lr{KNN} بود که توانایی کار با مجموعه‌ای
شامل ویژگی‌های عددی و دسته‌ای را دارد.

در مرحلهٔ پیش‌پردازش، داده‌ها برای استفاده در الگوریتم‌های مبتنی بر فاصله آماده‌سازی شدند. مقادیر
گمشده برای ویژگی‌های عددی با میانگین میانه و برای ویژگی‌های دسته‌ای با مد جایگزین شدند. سپس
ویژگی‌های عددی با استفاده از روش مقیاس‌گذاری \lr{Min--Max} نرمال‌سازی شدند تا از غلبهٔ ویژگی‌هایی با
دامنهٔ عددی بزرگ‌تر در محاسبهٔ فاصله جلوگیری شود. در نهایت، مجموعه‌داده با نسبت ۸۰/۲۰ به آموزش و
آزمون تقسیم شد تا ارزیابی مدل روی داده‌های دیده‌نشده به‌درستی انجام گیرد.

نتایج تحلیل اکتشافی الگوهای بالینی قابل‌توجهی را آشکار کردند. در میان تمام ویژگی‌های عددی،
\lr{thalach} (حداکثر ضربان قلب) و \lr{oldpeak} (افت قطعهٔ \lr{ST} ناشی از ورزش) بیشترین ارتباط را
با وجود بیماری قلبی نشان دادند. بیماران مبتلا به بیماری قلبی معمولاً مقدار \lr{thalach} پایین‌تر و
مقدار \lr{oldpeak} بالاتری داشتند و این موضوع منجر به تفکیک‌پذیری بهتر میان دو کلاس در نمودارهای
پراکندگی و نمودارهای جفتی شد. در مقابل، ویژگی‌هایی مانند \lr{chol} و \lr{trestbps} ارتباط‌های ضعیف‌تر
و پراکنده‌تری داشتند و شامل چند مقدار پرت بالینی قابل‌قبول بودند. همچنین، توزیع نسبتاً متعادل متغیر
هدف شرایط مناسبی را برای یادگیری مدل بدون نیاز به روش‌های متوازن‌سازی فراهم کرد.

طبقه‌بند \lr{KNN} با بهره‌گیری از یک تابع فاصلهٔ ترکیبی اقلیدسی–همینگ به‌صورت کامل پیاده‌سازی شد.
عملکرد مدل برای مقادیر مختلف $k$ مورد ارزیابی قرار گرفت و بهترین نتایج در مقدار $k = 3$ حاصل شد.
در این مقدار، مدل بالاترین دقت، دقت مثبت، بازخوانی و امتیاز \lr{F1} را به دست آورد. این نتیجه با الگوهای
یافت‌شده در تحلیل اکتشافی هم‌خوان است؛ زیرا ساختار داده‌ها به‌واسطهٔ چند ویژگی بسیار تفکیک‌کننده
ساختاری محلی دارند و مقدارهای کوچک‌تر $k$ بهتر می‌توانند این الگوهای محلی را بدون بیش‌ازحد هموار
کردن مرزها ثبت کنند. ماتریس اغتشاش نیز نشان داد که مدل توانایی خوبی در تشخیص صحیح هر دو کلاس دارد و
بیشتر موارد بیماری قلبی به‌درستی شناسایی شده‌اند.

به‌طور کلی، این مطالعه نشان می‌دهد که الگوریتم \lr{KNN} — با وجود سادگی آن — در صورت انجام
پیش‌پردازش مناسب و انتخاب تابع فاصلهٔ صحیح، می‌تواند عملکرد پیش‌بینی‌کنندهٔ بسیار خوبی روی داده‌های
پزشکی داشته باشد. نتایج همچنین اهمیت ویژگی‌های مرتبط با عملکرد ورزشی مانند \lr{thalach} و
\lr{oldpeak} را به‌عنوان شاخص‌های کلیدی بیماری قلبی برجسته می‌کند و نشان می‌دهد که ویژگی‌هایی
همچون \lr{chol} و \lr{trestbps} به‌تنهایی قدرت محدودتری در تمایز بین کلاس‌ها دارند.

برای آینده، پیشنهاد می‌شود از روش‌هایی مانند اعتبارسنجی متقابل (\lr{Cross-Validation}) برای برآورد
پایدارتر عملکرد مدل استفاده شود، وزن‌دهی بین فاصلهٔ عددی و دسته‌ای بهینه‌سازی گردد یا حتی مدل با
طبقه‌بندهای پیشرفته‌تر مانند رگرسیون لجستیک، درخت تصمیم، \lr{SVM} یا روش‌های ensemble مقایسه شود.
چنین توسعه‌هایی می‌توانند به بهبود دقت طبقه‌بندی و ارائهٔ بینش‌های عمیق‌تر دربارهٔ روابط میان
ویژگی‌ها در پیش‌بینی ریسک بیماری قلبی کمک کنند.